\section{Elección del grupo de comparación}
\label{anexo:control}

Este anexo desarrolla la comparación entre las tres definiciones posibles del
grupo de control que la subsección de Identificación resume. La conclusión es
que las nunca tratadas —el criterio de \textcite{FISCHER2025}— no es superior en
todos los frentes, sino la única definición que no reprueba de manera marcada en
ninguno.

Las tres definiciones son: el \emph{anillo}, formado por las incumbentes
situadas entre 2 y 5 kilómetros del evento; el grupo \emph{lejano}, formado por
las estaciones que nunca estuvieron a menos de 5 kilómetros de una entrada; y
las \emph{nunca tratadas}, que es la unión de las dos anteriores y equivale al
conjunto de estaciones que nunca reciben una entrada dentro del radio de
tratamiento. Las tres se estiman sobre la misma muestra de tratadas y con la
misma especificación, de modo que la única diferencia es contra quién se las
compara.

Conviene tener presente desde el inicio que las nunca tratadas no son un cuarto
grupo independiente sino la suma de los otros dos. Cualquier característica en
que aparezcan como intermedias lo son por construcción, al ser un promedio
ponderado de un grupo próximo a las tratadas y otro alejado, y no por una
propiedad sustantiva propia.

\subsection{Magnitud y precisión}

La \Cref{tab:control} del cuerpo muestra que las tres definiciones entregan
estimaciones muy próximas entre sí. Ese hecho es, por sí mismo, el resultado más
importante de este anexo: la conclusión no depende de la elección, de modo que
esta no es un grado de libertad del que dependa el signo o el orden de magnitud
del efecto.

Donde sí difieren es en precisión. Las nunca tratadas más que duplican el tamaño
del anillo y superan al grupo lejano, y en consecuencia entregan el error
estándar más bajo en los cuatro combustibles cuando el resultado es el precio, y
en dos de los tres cuando es el margen. La \Cref{tab:balance} reporta el número
de estaciones de cada grupo junto con sus características promedio antes del
tratamiento.

\begin{table}[H]
\centering
\caption{Balance de covariables y tamaño de cada grupo de comparación}
\label{tab:balance}
\tabla{tab_balance_control}
\end{table}

\subsection{Comparabilidad y contaminación}

La \Cref{tab:balance} deja ver por qué ninguno de los dos subconjuntos domina al
otro.

El anillo es el más parecido a las tratadas en cuatro de las cinco dimensiones
observadas. Comparte con ellas el entorno urbano y opera en mercados de densidad
comparable: 6,2 competidores dentro de 2 kilómetros contra 8,9 de las tratadas,
mientras que las lejanas apenas alcanzan 2,1. La brecha de margen es de 9,9
pesos por litro, frente a 38,8 del grupo lejano. Pero el anillo es también el
grupo contaminado: el ejercicio de atenuación por distancia del cuerpo muestra
que el efecto de la entrada sigue siendo distinto de cero en la banda de 2 a 3
kilómetros, que es precisamente su borde interior. Un control parcialmente
tratado atenúa la estimación hacia cero.

El grupo lejano es la imagen inversa. Está limpio de contaminación por
construcción, ya que ninguna de sus estaciones estuvo jamás a menos de 5
kilómetros de una entrada, pero es el peor balanceado en todas las dimensiones
salvo una. Sus mercados son estructuralmente más ralos, su margen promedio
supera en un $55\%$ al de las tratadas, y solo el $9\%$ de sus estaciones está
en la Región Metropolitana contra el $30\%$ de las tratadas. Usarlo en exclusiva
expone la estimación a que una diferencia de tendencia entre mercados urbanos y
rurales se confunda con el efecto de la entrada.

La excepción es la composición regional, y ahí el orden se invierte. El anillo se
desvía veintidós puntos por encima de las tratadas en participación
metropolitana y el grupo lejano veinte puntos por debajo, de modo que al unirlos
ambos sesgos se cancelan y las nunca tratadas quedan a un punto de las tratadas.
Es la ilustración más clara del punto anterior: la ventaja es aritmética, no
sustantiva, pero no por ello menos real para el diseño.

En las demás dimensiones las nunca tratadas heredan una parte de cada problema y
ninguna en su forma extrema. Conviene notar que la contaminación que hereda del
anillo opera en una dirección conocida: atenúa el efecto hacia cero, de modo que
la estimación resultante debe leerse como una cota inferior en magnitud y no
como un valor que podría estar inflado.

\subsection{Tendencias previas}

La \Cref{tab:pretend} reporta, para cada grupo y combustible, dos contrastes
sobre el período previo a la entrada. El primero es el test conjunto de que todos
los coeficientes anteriores a $\tau = -1$ son nulos. El segundo es la pendiente
de la trayectoria previa, junto con la deriva que esa pendiente acumularía en
doce meses expresada como fracción del efecto estimado.

\begin{table}[H]
\centering
\caption{Contrastes de tendencias previas por grupo de comparación}
\label{tab:pretend}
\tabla{tab_pretend_control}
\end{table}

La distinción entre ambos contrastes no es un tecnicismo. El test conjunto
evalúa la hipótesis nula contra alternativas arbitrarias y tiene, por esa misma
razón, poca potencia frente a una deriva lineal suave, que es justamente la
forma que importaría. El contraste de pendiente es el que tiene potencia contra
esa alternativa, y la razón entre la deriva acumulada y el efecto estimado es la
que permite juzgar si una pendiente estadísticamente indistinguible de cero es,
aun así, lo bastante grande como para explicar una fracción relevante del
resultado.

Bajo el contraste de pendiente ninguna de las tres definiciones se aproxima a la
significancia convencional: los valores $p$ van de $0{,}58$ a $0{,}97$. Las
derivas acumuladas, en cambio, no son despreciables en todos los casos, y su
lectura exige atender al signo. En once de las doce combinaciones de grupo y
combustible la deriva previa es de signo contrario al efecto estimado, es decir,
los precios de las tratadas venían subiendo respecto de los controles antes de
la entrada. Una pre-tendencia de ese signo no puede generar el resultado: opera
en contra de encontrarlo, y hace que la estimación sea conservadora.

La magnitud sí distingue entre grupos. En el grupo lejano la deriva alcanza el
$71\%$ del efecto estimado en la gasolina de 95 octanos y el $48\%$ en la de 97,
mientras que en las nunca tratadas no supera el $38\%$ y en el anillo el $41\%$.
Es un argumento adicional, y del mismo signo que los anteriores, contra usar el
grupo lejano en exclusiva.

\subsection{Síntesis}

La \Cref{fig:comparacion} superpone las tres trayectorias estimadas sobre la
misma muestra de tratadas. Se distinguen apenas entre sí, tanto antes como
después de la entrada.

\begin{figure}[H]
\centering
\includegraphics[width=\textwidth]{es_control_comparacion.pdf}
\caption{Estudio de eventos bajo las tres definiciones del grupo de comparación}
\label{fig:comparacion}
\end{figure}

Los criterios se resumen así. En precisión gana el conjunto de las nunca
tratadas; en comparabilidad gana el anillo, salvo en composición regional; en
ausencia de contaminación gana el grupo lejano. Las nunca tratadas es la única
alternativa que no queda en último lugar en ninguno de los tres, y es además el
criterio que emplea \textcite{FISCHER2025} en un diseño de las mismas
características, lo que hace que la elección tenga precedente y no dependa del
juicio del autor.
